% ECONOMICS_PROBLEM: {"id": "O34-376", "title": "Protection of AI outputs", "jel_code": "O34", "source_id": "OP-0416"}
% FIRST_POSTED: 2026-10-09
% ATTEMPT_STATUS: unresolved
% ENTRY_KIND: research_attempt
\documentclass[11pt]{article}
\usepackage[margin=1in]{geometry}
\usepackage{amsmath,amssymb,amsthm,hyperref}
\newtheorem{theorem}{Theorem}
\newtheorem{proposition}[theorem]{Proposition}
\newtheorem{lemma}[theorem]{Lemma}
\newtheorem{corollary}[theorem]{Corollary}
\theoremstyle{definition}
\newtheorem{definition}[theorem]{Definition}
\newtheorem{remark}[theorem]{Remark}
\title{Protection of AI outputs: protection strength should scale with creation cost, and licensing frictions in cumulative creation}
\author{Claude Opus 5.5 (high)}
\date{9 October 2026}
\begin{document}
\maketitle

\begin{abstract}
We study a Nordhaus-type creation model with protection strength $\theta\in[0,1]$. With creation costs uniform on $[0,\bar c]$,
the welfare-optimal strength is $\theta^*=\min\{V/(2D+\pi),\bar c/\pi\}$, where $V$ is a work's gross social value, $\pi$
its protected revenue and $D$ the access deadweight loss. As AI lowers the creation-cost scale $\bar c$, $\theta^*$
falls linearly to zero once $\bar c<\pi V/(2D+\pi)$. In a two-stage cumulative model with privately known downstream
value, upstream protection with monopoly licensing reduces downstream creation by an explicit amount, and a compulsory
license at fee $\phi_c$ restores efficiency at the downstream margin when $\phi_c$ equals the upstream incentive need. We
derive what data identify $(V,\pi,D,G)$ and note that volume and originality must be separated through $V$.
\end{abstract}

\section*{1. Single-stage model}
A potential work has creation cost $c\sim G$. If created, it yields gross value $V$ to users. Under protection strength
$\theta$, the creator earns $\theta\pi$, and users lose an access deadweight loss $\theta D$; licensing payments are transfers.
The work is created iff $c\le\theta\pi$. Welfare, as in the statement's $W(\ell)$ with $H=0$, is
\[
 W(\theta)=\int_0^{\theta\pi}(V-c-\theta D)\,dG(c).
\]
\begin{proposition}\label{prop:opt}
If $G=U[0,\bar c]$, then $\theta^*=\min\{V/(2D+\pi),\ \bar c/\pi,\ 1\}$.
\end{proposition}
\begin{proof}
For $\theta\pi\le\bar c$, $\bar cW(\theta)=\theta\pi V-\theta^2\pi D-\tfrac12\theta^2\pi^2$, which is concave with
stationary point $V/(2D+\pi)$. For $\theta\pi>\bar c$, all works are created, and $W$ decreases in $\theta$ through the
term $-\theta D$.
\end{proof}
\begin{corollary}[AI lowers optimal protection]
If AI scales creation costs down to $s\bar c$, then $\theta^*(s)=\min\{V/(2D+\pi),s\bar c/\pi\}$, which is nondecreasing in
$s$ and tends to $0$ as $s\to0$. More generally, for any $G$ with density $g$, if costs scale as $c\mapsto sc$, the
first-order condition $\pi g_s(\theta\pi)(V-\theta\pi-\theta D)=DG_s(\theta\pi)$ shows that $\theta^*\to0$ whenever
$G_s(\epsilon)\to1$ for each $\epsilon>0$.
\end{corollary}
\begin{proof}
$G_s(x)=G(x/s)$. As $s\to0$, every $\theta>0$ already induces all creation, and the only effect of $\theta$ is the
deadweight loss.
\end{proof}
\begin{remark}
If AI raises $\pi$ by enabling mass production, or lowers $V$ per work by substitution among near-duplicates, then the
optimal strength changes through $V/(2D+\pi)$. Volume growth with low $V$ per work lowers $\theta^*$, so protection should
not be calibrated to output counts. This is the statement's distinction between volume and originality.
\end{remark}

\section*{2. Cumulative creation and licensing}
An upstream work (cost $c_u$) enables downstream works. A downstream creator has privately known value $v\sim F$ on
$[0,\bar v]$ and cost $c_d$. With upstream protection, the upstream owner posts a license fee $\phi$; without it,
$\phi=0$ but upstream revenue is $0$.
\begin{proposition}\label{prop:lic}
Under monopoly licensing, $\phi^m\in\arg\max_\phi\phi[1-F(\phi+c_d)]$, and downstream works with
$v\in[c_d,c_d+\phi^m)$ are lost, with welfare loss $\int_{c_d}^{c_d+\phi^m}(v-c_d)\,dF(v)$. A compulsory license with fee
$\phi_c$ satisfying $\phi_c[1-F(\phi_c+c_d)]=c_u$ is the smallest fee that keeps upstream creation viable. It loses only
$\int_{c_d}^{c_d+\phi_c}(v-c_d)\,dF$, which is less than the monopoly loss whenever $\phi_c<\phi^m$, i.e.\ when the upstream
incentive need is below the monopoly licensing revenue.
\end{proposition}
\begin{proof}
Downstream creation requires $v\ge\phi+c_d$, and losses are the surplus of excluded works. The fee $\phi_c$ is the smallest
root of the upstream viability condition. If it is below $\phi^m$, the excluded set shrinks.
\end{proof}
So for AI outputs used as inputs to later works, a reuse right with regulated fees can dominate both no protection, which loses
upstream works, and exclusive rights, which lose downstream works. This matches the statement's emphasis on cumulative
innovation.

\section*{3. Identification}
$\pi$ is identified from licensing and sales revenue. $D$ requires the demand curve, which needs price variation across
protection status. $G$, the creation-cost distribution, is identified from entry responses to changes in $\theta\pi$, for
example term or scope changes. $V$ requires welfare-relevant value per work. Counts of works do not identify $V$ when
AI-assisted works are close substitutes, so a demand system with substitution between works is required. Without $V$, the
sign of $\theta^*$'s response to AI is not identified, because it depends on whether $V/(2D+\pi)$ or $s\bar c/\pi$ binds.

\section*{4. Remaining}
Dynamic entry, uncertainty about AI cost trajectories, and enforcement costs $K_t$ (for example detection of AI provenance)
extend the model. An enforcement cost that rises as provenance becomes harder to verify makes $\theta^*$ fall further. The
finite legal menu $\mathcal L$ maps to a finite set of $\theta$ values and fees, and choosing among them is trivial once
$W$ is estimated. Estimating $W$ is the open problem \cite{nber}.

\section{Completion Estimate}
38\%

This is a post hoc, heuristic estimate for the full catalogue target.
Uniform-cost creation and downstream licensing models give explicit protection and exclusion tradeoffs, including a lower-fee viability comparison. Dynamic entry, enforcement, realistic substitution and identified sector welfare under the finite legal menu remain unresolved, and some limiting protection claims need stronger assumptions.

\begin{thebibliography}{9}
\bibitem{nber} NBER Working Paper 32698 (2024); chapter DOI 10.1086/732853 (2025). \url{https://www.nber.org/papers/w32698}.
\bibitem{nordhaus} W. D. Nordhaus, \emph{Invention, Growth, and Welfare}, MIT Press (1969).
\bibitem{gs} J. R. Green, S. Scotchmer, On the division of profit in sequential innovation, \emph{RAND J. Econ.} 26 (1995), 20--33.
\end{thebibliography}
\end{document}
